\section{Additional Definitions}
\label{app:defs}

This appendix collects definitions used by the main text.

\subsection{Centering map and carries}
Let \(\BoundB\in\Z_{>0}\) and \(\Delta=2\BoundB+1\). Define the centering map
\(\centerfunc:[-2\BoundB,2\BoundB]\to[-\BoundB,\BoundB]\) by
\[
\centerfunc(x)=
\begin{cases}
x & \text{if } x\in[-\BoundB,\BoundB],\\
x+\Delta & \text{if } x\in[-2\BoundB,-\BoundB-1],\\
x-\Delta & \text{if } x\in[\BoundB+1,2\BoundB].
\end{cases}
\]
Applied coefficient-wise, this map allows unbiased combination of two bounded values when at least one is uniform.

\begin{lemma}[Centering preserves uniformity]\label{lem:center-uniform}
Fix any \(a\in[-\BoundB,\BoundB]\). If \(b\getsunif [-\BoundB,\BoundB]\) then \(\centerfunc(a+b)\) is uniform over
\([-\BoundB,\BoundB]\).
\end{lemma}

\subsection{Vanishing polynomials for membership constraints}
Membership constraints are expressed using vanishing polynomials over \(\Fq\) with signed lifting.

\begin{definition}[Vanishing polynomial for \(\lbrack -\BoundB,\BoundB\rbrack\)]\label{def:vanish-bound}
Let \(\BoundB\in\Z_{>0}\). Define
\[
  P_{\BoundB}(X) \;=\; \prod_{i=-\BoundB}^{\BoundB}\bigl(X-\langle i\rangle_q\bigr)\ \in\ \Fq[X].
\]
For \(x\in\Fq\), we have \(x\in[-\BoundB,\BoundB]\) (via signed lifting) iff \(P_{\BoundB}(x)=0\) in \(\Fq\).
\end{definition}

\begin{definition}[Vanishing polynomial for \(\{-1,0,1\}\)]\label{def:vanish-carry}
Define \(P_{1}(X)=(X+1)X(X-1)\in\Fq[X]\). Then \(P_1(x)=0\) iff \(x\in\{-1,0,1\}\).
\end{definition}
